% Ad Familiares 10.18 --- headnote
% ad-familiares-10-18, 18 May 43 BC, in camp on the march from the Is\`ere to Forum Voconii

\textit{L. Munatius Plancus to Cicero, written in camp on the
march from the Is\`ere down toward Forum Voconii on 18 May
43 BC --- Perseus dateline \textit{Scr. in castris ex itinere
ab Isara Forum Voconi xv K. Iun. a. 711 (43)}. This is the
dispatch in which Plancus explains, while it is still
happening, why he has decided to march south to join Lepidus
rather than wait, as caution would have counselled, for D.
Brutus to cross the Is\`ere with his army and come up to him.}

\textit{The strategic problem is acute. Antony has reached
Forum Iulii; Lepidus, with seven legions, is at Forum
Voconii twenty-four miles inland; Lepidus's army is
notoriously of doubtful loyalty (the ``inconstancy and
unreliability'' Plancus and Cicero both fear); Lepidus
himself is now writing pleading letters, doubled up by
even more urgent ones from his legate Juventius Laterensis.
If Plancus stays back on the safer line and Lepidus is
overwhelmed --- or, worse, peeled away by his own troops
--- the resulting collapse will be laid to Plancus's
account either as obstinacy or as cowardice. Better to
take the risk, push down to join Lepidus, and try to hold
the army together by physical presence. The fourth section
gives the operational details: he has broken camp on the
18th of May, but left the bridge over the Is\`ere standing,
with two forts and strong garrisons at its heads, so that
when D. Brutus comes up he can cross at once.}

\textit{The military-administrative register is at its
strongest here, with Plancus weighing alternatives in the
unmistakable voice of a serving commander explaining a
field decision to his political patron. The image of the
hidden wound is the sentence Shackleton Bailey marked as
Plancus's most arresting --- ``I cannot help shuddering,
if some wound is festering beneath the skin which can do
its damage before it can be either known or treated.''
And the contempt for Antony's quartermaster Ventidius
Bassus (``Ventidi mulionis castra'' --- ``the muleteer
Ventidius's camp,'' from the well-worn jibe that he had
once driven mule-trains in his youth) is the one piece
of openly partisan colour. Within ten days of writing this
letter, Plancus will have joined Lepidus, and the army
about whose loyalty he and Laterensis were so anxious
will have gone over to Antony en masse.}
